\documentclass[11pt, a4paper]{article}
\usepackage[utf8]{inputenc}
\usepackage[T1]{fontenc}
\usepackage{geometry}
\usepackage{amsmath, amssymb, amsfonts, amsthm}
\usepackage{parskip}
\usepackage{hyperref}
\usepackage{enumitem}
\usepackage[title, titletoc]{appendix}
\usepackage[super]{cite}
\usepackage[normalem]{ulem}     % if you want \uline (optional)
\usepackage[dvipsnames]{xcolor}   % ← add this line

\usepackage{hyperref}
\usepackage{xurl}               % <-- Highly recommended!

% Good modern setup
\hypersetup{
    colorlinks=true,
    linkcolor=blue,
    citecolor=blue,
    urlcolor=blue!70!black,
    pdfborder={0 0 0},          % no boxes around links
}

% Geometry settings
\geometry{left=1in, right=1in, top=1in, bottom=1in}

% Theorem environments
\newtheorem{theorem}{Theorem}[section]
\newtheorem{lemma}[theorem]{Lemma}
\newtheorem{proposition}[theorem]{Proposition}
\newtheorem{corollary}[theorem]{Corollary}
\newtheorem{definition}[theorem]{Definition}
\newtheorem{conjecture}[theorem]{Conjecture}
\newtheorem{axiom}{Axiom}

\theoremstyle{remark}
\newtheorem{remark}[theorem]{Remark}

% Custom commands
\newcommand{\Id}{\mathsf{Id}}
\newcommand{\Mem}{\mathsf{Mem}}
\newcommand{\RI}{\mathcal{I}}
\newcommand{\ZF}{\mathsf{ZF}}



\title{The Universality of Thomas L. Etter's Three-Place Identity:\\A Formal Proof and Its Significance}
\author{James Bowery}
\date{2025-12-24}

\begin{document}

\maketitle

\begin{abstract}
    We present a formal proof, mechanized in Lean 4, demonstrating that three-place relative identity\cite{quine} is expressively adequate for mathematics. Given any model of ZFC, we construct definitional bridges showing that membership can be derived from relative identity and vice versa, with the derived membership extensionally equivalent to the original.
\begin{tabular}{|c|c|}
\hline
\textbf{ZFC} & \textbf{Etter} \\
\hline
$=$ primitive & $x = y \;:=\; (\forall z,\, z \in x \leftrightarrow z \in y) \land (\forall w,\, x \in w \leftrightarrow y \in w)$ \\
\hline
$\in$ primitive & $y \in x \;:=\; \neg\,\mathrm{Id}(x,y,x)$ \\
\hline
(none) & $Quine \mathrm{Id}(x,y,z)$ primitive \\
\hline
\end{tabular} \\
    This establishes that ZFC can be faithfully reformulated using relative identity as the sole primitive relation. We discuss the philosophical significance of this result for foundations of mathematics.
\end{abstract}

\section{Introduction}

The question of foundational primitives in mathematics has occupied logicians since Frege. The standard answer---that set membership $\in$ suffices as the sole primitive, with all mathematical objects reducible to sets---has proven technically successful but philosophically unsatisfying. As Benacerraf\cite{benacerraf} observed, the arbitrariness of set-theoretic encodings (why is $2 = \{\emptyset, \{\emptyset\}\}$ rather than $\{\{\emptyset\}\}$?) suggests that sets may not capture the essence of mathematical objects. Tom Etter\cite{etter_eq,etter_notes,etter_mem,etter_3pl} proposed a radical alternative: that \emph{identity}, not membership, is the true primitive. But not the two-place identity $x = y$ of classical logic, which Quine showed to be expressively impoverished. Rather, a \emph{three-place} relative identity $x(y = z)$, read ``$x$ regards $y$ as the same as $z$.''

This paper presents a machine-verified proof that Etter's identity theory is \emph{universal}---"capable of expressing all of mathematics". The proof has been formalized in Lean 4.  The Lean4 source is available at \url{http://github.com/jabowery/RelativeIdentity}

\emph{Terminological clarification}

When we say that three-place relative identity is “sufficient to express all of mathematics,” we mean the axioms of ZFC can be formulated and faithfully interpreted using identity as the sole primitive, via definitional translations. This is merely relative interpretability in the usual logical sense. We do \emph{not} claim that the axioms of relative identity alone suffice to derive or model ZFC. The present result establishes \emph{expressive} adequacy and interpretability, not \emph{foundational} sufficiency in the model-theoretic sense.

\section{Three-Place Relative Identity}

\begin{definition}[Relative Identity Structure]
A \emph{relative identity structure} on a type $U$ consists of a ternary predicate $\Id : U \to U \to U \to \mathsf{Prop}$ satisfying:
\begin{enumerate}[label=(\roman*)]
    \item \textbf{Reflexivity:} $\forall x\, y.\; \Id(x, y, y)$
    \item \textbf{Symmetry:} $\forall x\, y\, z.\; \Id(x, y, z) \to \Id(x, z, y)$
    \item \textbf{Transitivity:} $\forall x\, y\, z\, w.\; \Id(x, y, z) \land \Id(x, z, w) \to \Id(x, y, w)$
\end{enumerate}
\end{definition}

The axioms assert that for each fixed ``viewpoint'' $x$, the relation $\Id(x, -, -)$ is an equivalence relation on the remaining arguments. This is weaker than requiring a single global equivalence: different viewpoints may identify different things.

\begin{remark}
The philosophical reading is: $x(y = z)$ means ``from $x$'s perspective, $y$ and $z$ are indistinguishable in the way that matters.'' This generalizes Quine's insight that identity within a formal language is indistinguishability with respect to that language's predicates.
\end{remark}

\section{The Two Bridges}

The proof of universality rests on two definitional bridges between identity and membership.

\begin{definition}[D1: Membership from Identity]
Given a relative identity structure $\RI$, define:
\[
y \in' x \;\stackrel{\text{def}}{=}\; \neg\, \Id(x, y, x)
\]
That is, $y$ exists for $x$ (is a member of $x$) when $x$ regards $y$ as \emph{different from itself}.
\end{definition}

The philosophical content is striking: existence is derived from distinguishability from non-existence. The viewpoint $x$ serves as a kind of ``ontological origin''---what exists for $x$ is precisely what $x$ distinguishes from the undifferentiated background (represented by $x$ itself).

\begin{definition}[D2: Identity from Membership]
Given a membership relation $\in$, define:
\[
\Id_\in(x, y, z) \;\stackrel{\text{def}}{=}\; (y \in x \land z \in x) \lor (\neg y \in x \land \neg z \in x)
\]
That is, $x$ regards $y$ and $z$ as identical when they have the same membership status with respect to $x$.
\end{definition}

\section{The Main Theorems}

\begin{theorem}[D2 Preserves Identity Structure]
\label{thm:d2-preserves}
For any membership relation $\in$ on $U$, the induced $\Id_\in$ satisfies the relative identity axioms.
\end{theorem}

\begin{proof}
We verify each axiom:

\textbf{Reflexivity:} For any $x, y$, either $y \in x$ or $\neg y \in x$. In the first case, $y \in x \land y \in x$ holds; in the second, $\neg y \in x \land \neg y \in x$ holds. Either way, $\Id_\in(x, y, y)$.

\textbf{Symmetry:} Immediate from the symmetry of $\land$ and $\lor$.

\textbf{Transitivity:} Suppose $\Id_\in(x, y, z)$ and $\Id_\in(x, z, w)$. We case-split:
\begin{itemize}
    \item If $y \in x \land z \in x$: From $\Id_\in(x, z, w)$ with $z \in x$, we cannot have $\neg z \in x$, so $z \in x \land w \in x$. Thus $y \in x \land w \in x$.
    \item If $\neg y \in x \land \neg z \in x$: From $\Id_\in(x, z, w)$ with $\neg z \in x$, we cannot have $z \in x$, so $\neg z \in x \land \neg w \in x$. Thus $\neg y \in x \land \neg w \in x$.
\end{itemize}
In both cases, $\Id_\in(x, y, w)$.
\qed
\end{proof}

\begin{theorem}[Round-Trip Equivalence]
\label{thm:roundtrip}
Let $\in$ be a membership relation satisfying the Irreflexivity of membership axiom ($\forall x.\, \neg x \in x$). Let $\RI = \Id_\in$ be the induced identity structure, and let $\in'$ be the membership derived from $\RI$ via D1. Then:
\[
\forall y\, x.\; y \in' x \iff y \in x
\]
\end{theorem}

\begin{proof}
Expanding definitions:
\begin{align*}
y \in' x &\iff \neg\, \Id_\in(x, y, x) \\
&\iff \neg\, [(y \in x \land x \in x) \lor (\neg y \in x \land \neg x \in x)]
\end{align*}

By Irreflexivity of membership, $\neg x \in x$. Thus $y \in x \land x \in x$ is false, and $\neg x \in x$ is true. The expression simplifies to:
\begin{align*}
&\iff \neg\, [\mathsf{false} \lor (\neg y \in x \land \mathsf{true})] \\
&\iff \neg\, (\neg y \in x) \\
&\iff y \in x
\end{align*}
\qed
\end{proof}

\begin{theorem}[Universality Theorem]
\label{thm:universality}
Identity theory is open-ended: the $\ZF$ axioms expressed in terms of derived membership $\in'$ are consistent with the relative identity axioms.
\end{theorem}

\begin{proof}
Let $M$ be a model of $\ZF$ with membership $\in_M$. By Theorem~\ref{thm:d2-preserves}, $\Id_{\in_M}$ is a relative identity structure. By Theorem~\ref{thm:roundtrip}, the derived membership $\in'$ is logically equivalent to $\in_M$. Therefore, any $\ZF$ axiom true of $\in_M$ is equally true of $\in'$. Since the $\ZF$ axioms are consistent (by assumption) and the identity axioms are satisfied by construction, their conjunction is consistent. Since $\ZF$ is a universal foundation for mathematics, identity theory augmented with the $\ZF$ axioms on derived membership can express all of mathematics.
\qed
\end{proof}

\section{Formalization in Lean 4}

The proof has been mechanized in Lean 4. The key definitions are:

\begin{verbatim}
structure RelativeIdentity (U : Type u) where
  Id : U → U → U → Prop
  refl : ∀ x y : U, Id x y y
  symm : ∀ x y z : U, Id x y z → Id x z y
  trans : ∀ x y z w : U, Id x y z → Id x z w → Id x y w

def MemFromId (RI : RelativeIdentity U) (y x : U) : Prop :=
  ¬ RI.Id x y x

def IdFromMem (mem : U → U → Prop) (x y z : U) : Prop :=
  (mem y x ∧ mem z x) ∨ (¬ mem y x ∧ ¬ mem z x)
\end{verbatim}

The round-trip theorem is proved using the \texttt{tauto} tactic after unfolding definitions and introducing the Irreflexivity of membership hypothesis:

\begin{verbatim}
theorem ZFModel_roundtrip (M : ZFModel U) (y x : U) :
    Mem'_from_ZFModel M y x ↔ M.mem y x := by
  unfold Mem'_from_ZFModel MemFromId RelativeIdentity_from_ZFModel
  unfold IdFromMem_is_RelativeIdentity IdFromMem
  simp only
  have nxx : ¬ M.mem x x := M.foundation x
  tauto
\end{verbatim}

\section{Philosophical Significance}

\subsection{Identity Precedes Existence}

Gian-Carlo Rota's slogan ``Identity precedes existence''\cite{rota} receives rigorous content through Etter's construction. The traditional view takes existence as primitive: things first exist, then we ask whether they are identical. Etter inverts this: identity (from a viewpoint) is primitive, and existence is derived as distinguishability from non-existence. Definition D1 makes this precise: $y$ exists for $x$ exactly when $x$ distinguishes $y$ from itself-as-null-viewpoint. Existence is not an absolute property but a relational one---to exist is to exist \emph{for} some perspective.

\subsection{The Relativity of Ontology}

Classical set theory posits a single, absolute universe of sets. Etter's framework suggests a more perspectival ontology: what exists depends on the viewpoint. Different ``observers'' $x$ may recognize different entities as existing.

This resonates with developments in:
\begin{itemize}
    \item \textbf{Quantum mechanics:} The observer-dependence of measurement outcomes
    \item \textbf{Constructive mathematics:} Existence requiring a construction, not just non-contradiction
    \item \textbf{Category theory:} Objects defined only up to isomorphism, with identity structure-dependent
\end{itemize}

\subsection{Connection to Quine's Role Identities}

Quine showed that identity within a formal language $L$ can be defined as the conjunction of all ``role identities''---the indistinguishability relations induced by each predicate position. Etter's three-place identity generalizes this by making the language/viewpoint itself a variable. We proved:
\[
\Id_\in(x, y, z) \iff (y \in x \leftrightarrow z \in x)
\]
This shows that $\Id_\in$ is exactly the Quine identity for the predicate $\lambda w.\, w \in x$. The viewpoint $x$ determines which predicate is relevant for identity.

\subsection{Mathematics Without Sets?}

The Universality Theorem does not eliminate sets---it shows they can be \emph{derived} from identity. This is analogous to how Dedekind derived real numbers from rationals: the reals don't disappear, but their foundational status changes. The philosophical gain is conceptual economy: instead of two primitive notions (membership and identity), we need only one (relative identity). The loss of absolute identity may be a feature, not a bug, for those troubled by puzzles of identity through time, across possible worlds, or in quantum contexts.

\section{Coordinating Relative Identities}

Distinguishing identities and relating them are already familiar from standard logic. Predicate calculus distinguishes objects by the predicates they satisfy, and it relates those objects by logical structure among predicates. What it does \emph{not} provide is a principled way to \emph{coordinate} multiple, independent identity judgments into a single, coherent mathematical structure. Three-place relative identity supplies exactly this missing capability.

\subsection{Identity as an Indexed Equivalence Relation}

A relative identity structure assigns to each viewpoint $x$ an equivalence relation $\Id(x,-,-)$. Identity is therefore not a single global relation, but a family of equivalence relations indexed by objects of the domain itself. This indexing is crucial. It means that:
\begin{itemize}
    \item different viewpoints may induce different partitions of the same universe;
    \item these partitions coexist on a common underlying domain;
    \item identity judgments made from one viewpoint can be systematically compared with those made from another.
\end{itemize}

The act of \emph{coordinating identities} consists precisely in aligning these multiple equivalence relations so that structure can be transferred across viewpoints.

\subsection{Coordination and Mathematical Structure}

Mathematical structure arises not merely from classification, but from the ability to \emph{use multiple classifications together}. Examples familiar from standard mathematics include:
\begin{itemize}
    \item ordered pairs, which coordinate two independent dimensions of identity;
    \item relations, which coordinate identities across domains;
    \item functions, which coordinate identities in a source with identities in a target;
    \item isomorphism, which coordinates identity judgments across different structures.
\end{itemize}

None of these constructions are possible with a single equivalence relation alone. They require multiple, independently meaningful notions of ``sameness'' that can be brought into alignment. Relative identity provides this alignment intrinsically: different viewpoints $x$ act as selectors of identity criteria, and mathematics proceeds by coordinating the resulting equivalence structures.

\subsection{Stereo Equality as a Discrete Case}

The Appendix makes this especially clear in the case of \emph{stereo equality}. Two independent equivalence relations $E_1$ and $E_2$ already suffice to generate structure that a single equivalence relation cannot: ordered pairs, relations, and membership-like behavior. Stereo equality may be viewed as a discrete approximation to three-place identity:
\begin{itemize}
    \item a single equality is monocular;
    \item two equalities provide binocular depth;
    \item three-place identity allows a continuous family of perspectives.
\end{itemize}

What matters is not the number three \emph{per se}, but the availability of multiple, coordinatable identity dimensions.

\subsection{Coordination Replaces Primitive Membership}

From this perspective, set membership is not fundamental. Membership is simply one way of coordinating identity judgments: an object $y$ is a member of $x$ exactly when $x$ distinguishes $y$ from itself under its identity criterion. Thus, mathematics does not require a primitive notion of ``being an element of.'' It requires only enough structure to:
\begin{enumerate}
    \item distinguish identities,
    \item relate them logically,
    \item coordinate them across perspectives.
\end{enumerate}

Three-place relative identity supplies all three.

\section{Future Directions}

Several extensions suggest themselves:

\begin{enumerate}
    \item \textbf{Non-well-founded variants:} The proof uses Irreflexivity of membership essentially. What happens with non-well-founded set theories (Aczel's AFA)?
    \item \textbf{Categorical semantics:} Three-place identity suggests a kind of ``indexed equivalence relation.'' This may connect to fibered categories or dependent type theory.
    \item \textbf{Quantum logic:} Etter hints at connections to von Neumann's quantum logic. Formalizing this could yield new insights into the foundations of quantum mechanics.
\end{enumerate}

\section{Conclusion}

We have presented a machine-verified proof that three-place relative identity is foundationally universal. The key insight is the mutual interpretability between identity and membership through two simple definitions. The philosophical upshot is that identity---understood as perspective-relative indistinguishability---may be a more fundamental notion than existence or membership. As Etter wrote: ``The statement `$x(y=z)$' is in effect a third-person statement... in order to understand `$x$ regards $y$ as the same as $z$' I must understand what it means for \emph{me} to think `$y$ and $z$ are the same' and then put myself in $x$'s shoes.'' Identity theory brings the subject into the foundations of mathematics.

\section*{Acknowledgments}

This formalization builds on Tom Etter's never-published papers, primarily ``Three-place Identity'' and ``Equalities and Quine Identities,'' developed at the Boundary Institute. The Lean formalization uses the Mathlib library\cite{mathlib}. Coding assistance and reviews were provided by the least expensive tier of the 4 major large language modes available during late 2025 (in no particular order): ChatGPT, Claude, Gemini and Grok.

\begin{thebibliography}{9}

\bibitem{quine}
Quine WVO. The Scope of Logic. In: Philosophy of Logic. Harvard University Press; 1986.

\bibitem{benacerraf}
Benacerraf P. What Numbers Could Not Be. Philosophical Review. 1965;74:47–73.

\bibitem{etter_eq}
    Etter T. Equalities and Quine Identities, Sections 1--6 [Internet; received and attested by author]. Boundary Institute; 2001. Available from \url{https://groups.io/g/lawsofform/files/Boundary%20Institute/Tom%20Etter%20Papers/Equality%20and%20Quine%20IDs%20s1-6.pdf}

\bibitem{etter_notes}
    Etter T. Notes on Identity and Pairing. [Internet; received and attested by author]; 2001. Available from \url{https://groups.io/g/lawsofform/files/Unpublished%20Tom%20Etter%20Papers%20Not%20Provided%20by%20Boundary%20Institute/Notes%20on%20Identity%20and%20Pairing.pdf}

\bibitem{etter_mem}
    Etter T. Membership and Identity [Internet]. Boundary Institute; 2006. Available from: \url{https://web.archive.org/web/20150925072222/http://www.boundaryinstitute.org/bi/articles/Membership_&_Identity.pdf}

\bibitem{etter_3pl}
    Etter T. Three-place Identity [Internet]. Boundary Institute; 2006. Available from: \url{https://web.archive.org/web/20150925072222/http://www.boundaryinstitute.org/bi/articles/Three-place_Identity.pdf}

\bibitem{etter_power}
    Etter T. The Expressive Power of Equality [Internet; received and attested by author]. Boundary Institute; 20 June 2001. (Contains the cell construction and the proof of the RCV$\to$ZF expressiveness theorem.)

\bibitem{rota}
Rota GC. The Primacy of Identity. In: Indiscrete Thoughts. Birkh\"auser; 1997.

\bibitem{mathlib}
    The Mathlib Community. Mathlib: The Lean Mathematical Library [Internet]. Available from: \url{https://github.com/leanprover-community/mathlib4}

\end{thebibliography}

\clearpage
\begin{appendices}
\section{The Stereo Equality Conjecture}
\label{app:stereo}

Alongside the three-place identity of the main body, Etter proposed a discrete simplification he called the \emph{Stereo Equality Theorem}: that \emph{two} equality predicates, taken together, already have the expressive power of ZF set theory. A single equality predicate is expressively dead-ended --- it distinguishes only equivalence classes --- whereas two predicates provide the ``binocular depth'' needed to encode order, and hence pairs, relations, and membership.

At present this is a \emph{conjecture}. Etter announces it in \emph{Three-place Identity}\cite{etter_3pl} and defers the proof to Section~3 of \emph{Membership and Identity}\cite{etter_mem}; the surviving draft of the latter develops the relevant machinery (intrinsic identity, \emph{division}, the category of \emph{Other}) but breaks off before reaching the theorem. Our contribution is to state the conjecture precisely enough to be proved or refuted, and to record it as a machine-checkable proposition in Lean.

\subsection{Why a symmetric combination cannot suffice}

Write $E_1$ and $E_2$ for the two equalities, each reflexive, symmetric and transitive. The simplest candidate for a derived membership is the intersection
\[
y \in' x \iff \neg\bigl(E_1(y, x) \land E_2(y, x)\bigr).
\]
Because $E_1$ and $E_2$ are symmetric, this relation is symmetric: $y \in' x \iff x \in' y$. Set membership is not symmetric --- $\emptyset \in \{\emptyset\}$ while $\{\emptyset\} \notin \emptyset$ --- so no symmetric combination of the two equalities can be a membership relation. The \emph{directedness} of membership must instead be supplied by a directed \emph{link} between the two dimensions. (This obstruction is recorded in Lean as \texttt{Stereo.Refutation}.)

\subsection{Etter's link construction}

The construction that supplies this directedness is sketched in \cite{etter_3pl}, footnote~5. It is built not on two transitive equalities but on a single equality \emph{with transitivity dropped}.

\begin{definition}[Weak equality]
$[x{=}y]$ is a predicate that is reflexive and symmetric but \emph{not} necessarily transitive.
\end{definition}

\begin{definition}[Meet]
$[x =: y, z, \dots]$ holds when $x$ equals each of $y, z, \dots$ and anything unequal to $y$ is unequal to $x$, ditto $z$: that is, $x$ lies at the meet of their neighbourhoods.
\end{definition}

\begin{definition}[Link membership]
\[
x \in' y \;:=\; \exists\, x',x'',y',y'',q,q',q''.\;
[q{=:}x,q'] \;\&\; [q'{=:}y,q,q''] \;\&\; [x'{=:}x,x''] \;\&\; [x''{=:}x,x'] \;\&\; [y'{=:}y,y''] \;\&\; [y''{=:}y,y'].
\]
\end{definition}

\begin{conjecture}[Stereo Equality]
There is a model of the weak equality $[x{=}y]$ in which the derived link membership $x \in' y$ is the membership relation of a model of ZF set theory.
\end{conjecture}

These definitions are transcribed in \texttt{Stereo.lean} as \texttt{WeakEq}, \texttt{Meet}, and \texttt{MemPrime}, and the conjecture as the proposition \texttt{etterStereoConjecture}. It is stated deliberately as a \texttt{Prop} rather than a proved \texttt{theorem}: a future paper may establish \texttt{etterStereoConjecture}, or its negation, in Lean. Footnote~5 uses a single non-transitive equality, which Etter calls ``the heart'' of the two-equality theorem; isolating the precise two-equality packaging is itself part of the open problem.

\subsection{Relation to three-place identity}

Stereo equality is a discrete special case of three-place identity:
\begin{enumerate}
    \item \textbf{Monocular (1-place):} $x = y$, classical identity --- sees only equivalence classes.
    \item \textbf{Binocular (2-place):} $E_1(x, y)$ and $E_2(x, y)$ --- sees depth and structure.
    \item \textbf{Omnicular (3-place):} $x(y = z)$ --- the viewpoint $x$ varies.
\end{enumerate}
A three-place identity $x(y=z)$ restricts to a stereo equality at two fixed viewpoints $x_1, x_2$, via $E_1(y, z) := x_1(y=z)$ and $E_2(y, z) := x_2(y=z)$; conversely a stereo equality lifts to a three-place identity through a selector that toggles between $E_1$ and $E_2$.

\section{The RCV Language Theorem}
\label{app:rcv}

In ``Notes on Identity and Pairing,'' Etter introduces a canonical class of languages known as RCV languages. These languages serve as a universal basis for his identity theory, asserting that even complex logical systems can be reduced to three specific identity predicates.

\subsection{Definitions}

\begin{definition}[Slang]
A \emph{Slang} (Standard Language) is defined as the predicate calculus augmented with a finite list of primitive predicates, containing no terms other than variables.
\end{definition}

\begin{definition}[RCV Language]
An \emph{RCV Language} is a Slang $I$ whose only primitive predicates are three global identity predicates, denoted as:
\[
R(x=y), \quad C(x=y), \quad \text{and} \quad V(x=y)
\]
\end{definition}

\subsection{Theorems}

The significance of RCV languages lies in their universality and their ability to constructively generate set theory.

\begin{theorem}[RCV Reduction]
Any Slang $S$ is a list variant of an RCV language $I$. That is, any standard language can be transformed into an RCV language effectively by reordering predicates.
\end{theorem}

\begin{theorem}[Constructibility of ZF]
ZF set theory can be constructively defined in terms of an RCV language. This is Etter's expressiveness theorem, proved in full in \emph{The Expressive Power of Equality}\cite{etter_power} by the cell construction (below) and mechanized in \texttt{Cell.lean}.
\end{theorem}

\begin{theorem}[Equivalence to Pairing]
If a Slang $S$ has the expressive power to define a global pairing predicate, then one can effectively construct its corresponding RCV language $I$.
\end{theorem}

\subsection{The cell construction}

Etter's RCV$\to$ZF membership is the \emph{cell construction} of \emph{The Expressive Power of Equality}\cite{etter_power}. A \emph{cell} is an ordered triple $\langle x, y, n\rangle$ with $x \,\varepsilon\, y$ and $n \in \{1, 2\}$, whose value is $\mathrm{Val}(\langle x, y, n\rangle) = x$ if $n = 1$ and $y$ otherwise. The predicates $R$, $C$, $V$ are the row, column, and value equalities on cells (same $\langle x, y\rangle$, same $n$, same value), with $V$ taken as identity. Both ``first column'' and membership are recovered from $R, C, V$ alone:
\[ \mathrm{First}(c) \iff \forall d\,\exists e\,\bigl(V(e{=}d) \land C(e{=}c)\bigr), \]
\[ c\,E\,d \iff \exists e, f\,\bigl(R(e{=}f) \land \mathrm{First}(e) \land V(c{=}e) \land \neg\mathrm{First}(f) \land V(d{=}f)\bigr). \]
The argument-asymmetry carried by $\mathrm{First}/\neg\mathrm{First}$ is what makes $cEd$ a directed membership. This construction is mechanized in \texttt{Cell.lean} over Mathlib's \texttt{ZFSet}: Etter's Theorems~1 and~2 are \texttt{first\_iff} and \texttt{mem\_iff}, with a clean \texttt{\#print axioms}, and \texttt{RCV.lean} defines RCV membership as this $cEd$.

\subsection{Relation to stereo equality}

The RCV language is a ``trinocular'' extension of the stereo structure of Appendix~\ref{app:stereo}: stereo equality uses two equalities, RCV three. The three-equality result is the proved cell construction above; whether \emph{two} equalities already suffice is the open Stereo Equality Conjecture of Appendix~\ref{app:stereo}.

\end{appendices}
\end{document}
